\section{Stolen Dreams}

Have you ever dreamed of being a poppy? It is a life to envy. The rising sun warms your petals just as the breeze is keeping them cool. Occasionally, gentle rains wash everything clean. Then the bees arrive and tickle your stamen; there is no sweeter or pleasurable feeling. The anticipation is nearly unbearable.

But the poppy is unusual. Not content to be a plant, the poppy desires her own dreams. So the poppy spirit makes a deal. She will take away your pains, assuage your anxieties, eliminate your fears. In exchange, she uses your consciousness to dream her dreams as you passively watch them unfold.

The poppy becomes addicted to these dreams and will not let you go. She demands that you take her sweet nectar, no matter how difficult or how costly. And you? You would prefer to die while high rather than to wake up.

\paragraph{Recital}

\url{https://www.gornahoor.net/wp-content/uploads/2023/03/Poppy2.mp3}

\url{https://www.youtube.com/watch?v=qFLw26BjDZs}

\flrightit{Posted on 2023-03-01 by Cologero}

\begin{center}* * *\end{center}

\begin{footnotesize}
\begin{sffamily}
\texttt{Santiago on 2023-07-20 at 00:09 said:}

``The sacred soldiery, those who were once mortal and who were redeemed by Christ... .are seated upon the thrones of the Mystic Rose... The second soldiery is of the angels who never left heaven. They soar above the Rose like heavenly bees, in constant motion between the Rose and the radiance of God. Unlike earthly bees, however, it is from God, the mystical hive of grace, that they bring the sweetness to the flower, bearing back to God... the bliss of the souls of heaven.

The first host is more emphatically centered on the aspect of God as the Son; the second, on the aspect of God as the Father.''

from translator John Ciardi's notes to Canto 31, Paradiso

\hfill

\texttt{Aria on 2026-03-01 at 13:54 said:}

Charles... so pleased you recorded your voice. Three years ago today.

\hfill

\texttt{Juliano on 2026-03-24 at 13:08 said:}

Aria, thanks for commenting here. It brought my attention to this audio. I have wrote to Charles a few times, but never heard his voice before.


\hfill

\end{sffamily}
\end{footnotesize}

